\vspace{-8pt}
\section{Conclusion}
This paper set out to answer whether unintended drift across a chain of instruction-based image edits can be measured automatically, whether it compounds, and whether cheap mitigations reduce it. RQ1 and RQ3 answer yes, by a wide margin: region-locking cuts mean cumulative drift by 89.1\% and masked conditioning by 68.6\%, both at $p<0.0001$. RQ2 is less clean: drift appears front-loaded into occasional catastrophic failures rather than gradually compounding, a pattern the Drift Score as built cannot fully distinguish from no compounding at all. Two specific mechanisms, a measurement ceiling effect and a metric-mitigation margin mismatch, explain why two of the original hypotheses reversed rather than simply failing to replicate, evidence that the measurement captures something real rather than noise.

All results are specific to InstructPix2Pix, and the human-perceptual validation step was built but not administered; resolving both, together with testing against a second editor, is the next step toward a benchmark the field can rely on rather than a single-project case study.
